% Options for packages loaded elsewhere
\PassOptionsToPackage{unicode}{hyperref}
\PassOptionsToPackage{hyphens}{url}
\documentclass[
  ignorenonframetext,
]{beamer}
\newif\ifbibliography
\usepackage{pgfpages}
\setbeamertemplate{caption}[numbered]
\setbeamertemplate{caption label separator}{: }
\setbeamercolor{caption name}{fg=normal text.fg}
\beamertemplatenavigationsymbolsempty
% remove section numbering
\setbeamertemplate{part page}{
  \centering
  \begin{beamercolorbox}[sep=16pt,center]{part title}
    \usebeamerfont{part title}\insertpart\par
  \end{beamercolorbox}
}
\setbeamertemplate{section page}{
  \centering
  \begin{beamercolorbox}[sep=12pt,center]{section title}
    \usebeamerfont{section title}\insertsection\par
  \end{beamercolorbox}
}
\setbeamertemplate{subsection page}{
  \centering
  \begin{beamercolorbox}[sep=8pt,center]{subsection title}
    \usebeamerfont{subsection title}\insertsubsection\par
  \end{beamercolorbox}
}
% Prevent slide breaks in the middle of a paragraph
\widowpenalties 1 10000
\raggedbottom
\AtBeginPart{
  \frame{\partpage}
}
\AtBeginSection{
  \ifbibliography
  \else
    \frame{\sectionpage}
  \fi
}
\AtBeginSubsection{
  \frame{\subsectionpage}
}
\usepackage{iftex}
\ifPDFTeX
  \usepackage[T1]{fontenc}
  \usepackage[utf8]{inputenc}
  \usepackage{textcomp} % provide euro and other symbols
\else % if luatex or xetex
  \usepackage{unicode-math} % this also loads fontspec
  \defaultfontfeatures{Scale=MatchLowercase}
  \defaultfontfeatures[\rmfamily]{Ligatures=TeX,Scale=1}
\fi
\usepackage{lmodern}
\ifPDFTeX\else
  % xetex/luatex font selection
\fi
% Use upquote if available, for straight quotes in verbatim environments
\IfFileExists{upquote.sty}{\usepackage{upquote}}{}
\IfFileExists{microtype.sty}{% use microtype if available
  \usepackage[]{microtype}
  \UseMicrotypeSet[protrusion]{basicmath} % disable protrusion for tt fonts
}{}
\makeatletter
\@ifundefined{KOMAClassName}{% if non-KOMA class
  \IfFileExists{parskip.sty}{%
    \usepackage{parskip}
  }{% else
    \setlength{\parindent}{0pt}
    \setlength{\parskip}{6pt plus 2pt minus 1pt}}
}{% if KOMA class
  \KOMAoptions{parskip=half}}
\makeatother
\setlength{\emergencystretch}{3em} % prevent overfull lines
\providecommand{\tightlist}{%
  \setlength{\itemsep}{0pt}\setlength{\parskip}{0pt}}
\usepackage{bookmark}
\IfFileExists{xurl.sty}{\usepackage{xurl}}{} % add URL line breaks if available
\urlstyle{same}
\hypersetup{
  hidelinks,
  pdfcreator={LaTeX via pandoc}}

\author{\texorpdfstring{}{}}
\date{}

\begin{document}

\begin{frame}[fragile]{The AI Collaboration Model}
\protect\phantomsection\label{the-ai-collaboration-model}
The
\hyperlink{documentation-duality-and-the-agentic-skill-architecture}{Documentation
Duality} standard and three-tier skill architecture represent an
empirical bet: that structured, machine-readable documentation
measurably improves AI agent performance in research codebases. This
section reports our key observations.

The documentation investment creates a positive feedback loop: as agents
produce higher-quality outputs from structured context, developers
maintain that documentation, which in turn improves future interactions
{[}@lau2025aicoding{]}. We observed this concretely during
\texttt{template/} development---each module's \texttt{SKILL.md} was
refined through iterative AI-assisted generation, serving as both input
prompt and output validator.

The \texttt{SKILL.md} layer, with its MCP-aligned YAML frontmatter
{[}@anthropic2024mcp{]}, provides a critical bridge to the agentic
software paradigm. Lu et al.'s AI Scientist {[}@lu2024aiscientist{]}
demonstrates end-to-end autonomous research; Wang et al.'s OpenHands
{[}@wang2024opendevin{]} achieved 53\% on SWE-Bench Verified
{[}@jimenez2024swebench{]}---the first open-source system to exceed
50\%. These systems require structured, protocol-aligned tool
inventories to navigate unfamiliar codebases. An OpenHands-class agent
navigating \texttt{template/} reads \texttt{CLAUDE.md} for global
constraints, scans \texttt{AGENTS.md} for module surfaces, and invokes
capabilities via \texttt{SKILL.md} without hallucinating function
signatures. All 12 infrastructure modules carry \texttt{AGENTS.md} and
\texttt{README.md}; the 10 active subpackages additionally carry
\texttt{SKILL.md}, ensuring no documentation blind spots.

This three-tier model is, to our knowledge, novel in the research
software engineering literature. The scale of the
investment---approximately 170 documentation files across
\texttt{docs/}, \texttt{AGENTS.md}/\texttt{README.md} pairs, and
per-module skill descriptors---represents a deliberate commitment to
machine-readable context that reduces hallucination surface area.
\end{frame}

\begin{frame}[fragile]{The Learning Curve}
\protect\phantomsection\label{the-learning-curve}
The Thin Orchestrator pattern imposes a cognitive overhead on
researchers accustomed to writing monolithic scripts. The requirement to
factor all logic into \texttt{src/} modules and use scripts only as
stateless wiring introduces an additional layer of indirection. We
mitigate this through:

\begin{enumerate}
\tightlist
\item
  \textbf{Template examples}: The \texttt{code\_project} serves as a
  fully worked exemplar with comprehensive comments.
\item
  \textbf{Documentation Duality}: Every directory has both
  \texttt{README.md} (for humans) and \texttt{AGENTS.md} (for AI
  collaborators), reducing the cost of navigation.
\item
  \textbf{Interactive orchestrator}: \texttt{run.sh} provides a TUI menu
  that abstracts pipeline complexity.
\item
  \textbf{Skill-level documentation}: The \texttt{docs/guides/}
  directory provides four progressive guides (Levels 1--3 Beginner, 4--6
  Intermediate, 7--9 Advanced, 10--12 Expert) alongside a comprehensive
  new-project setup checklist.
\end{enumerate}
\end{frame}

\begin{frame}[fragile]{Limitations}
\protect\phantomsection\label{limitations}
\begin{itemize}
\tightlist
\item
  \textbf{LaTeX dependency}: The rendering pipeline requires a full TeX
  distribution (TeX Live or MiKTeX), which is a 4--6 GB install. This is
  the largest single dependency and is a barrier for researchers without
  system-level package management access.
\item
  \textbf{Python-centric}: The infrastructure layer is Python-only.
  Projects in other languages can use the rendering and steganography
  stages but cannot leverage the \texttt{scientific} or
  \texttt{validation} modules.
\item
  \textbf{Single-machine}: The pipeline runs locally. Distributed
  execution (e.g., across a compute cluster) is not natively supported,
  a gap where Snakemake, Nextflow, and CWL have clear superiority.
\item
  \textbf{Steganographic robustness}: Alpha-channel overlays are
  stripped by aggressive PDF optimization tools (e.g.,
  \texttt{qpdf\ -\/-optimize}). QR codes are visible and removable. The
  current system provides \emph{tamper detection} (via SHA-256 hashing)
  but not \emph{non-repudiation} in the cryptographic sense---it lacks
  private-key digital signatures. An attacker with access to the source
  code could reproduce the watermark without having run the original
  pipeline.
\item
  \textbf{Test duration}: The Zero-Mock policy increases test execution
  time from sub-second (mocked) to multi-minute (real) for the full
  infrastructure suite. This is acceptable for research workflows but
  may not suit continuous deployment scenarios.
\item
  \textbf{AI-native writing tools}: \texttt{template/} does not include
  an integrated AI writing assistant comparable to Overleaf's Copilot
  features or OpenAI Prism's GPT-5.2 context-aware editing. The
  \texttt{infrastructure.llm} module provides LLM review as a pipeline
  stage but not as an interactive writing environment.
\end{itemize}
\end{frame}

\end{document}
